%!TeX root=MemoriaTFG.tex

\chapter{Diseño experimental: modelos, protocolos y métricas}
\label{cap:experimental}

Este capítulo describe el conjunto de modelos evaluados en el
trabajo, los protocolos experimentales aplicados sobre los datasets
caracterizados en el capítulo~\ref{cap:datos} y las métricas
empleadas para reportar los resultados. La descripción se acota al
diseño y a los hiperparámetros: ninguna cifra de rendimiento aparece
en este capítulo. Los valores cuantitativos obtenidos se reservan al
capítulo~\ref{cap:resultados}.

El capítulo se estructura en cuatro bloques. El primero
(sección~\ref{sec:hardware-cap4}) describe el equipo físico y el
entorno de software sobre el que se ejecutan todos los
experimentos. El segundo (sección~\ref{sec:modelos-evaluados})
enumera los modelos evaluados, agrupados por familia y con sus
características de tamaño y entrada. El tercero
(secciones~\ref{sec:protocolos-base}--\ref{sec:protocolos-auxiliares})
describe los protocolos experimentales aplicados, sus
hiperparámetros y las decisiones metodológicas asociadas. El
cuarto (sección~\ref{sec:metricas-cap4}) formaliza las métricas
utilizadas y, por último, la sección~\ref{sec:reproducibilidad}
recoge los elementos de trazabilidad y reproducibilidad del
trabajo.

\section{Hardware y entorno de ejecución}
\label{sec:hardware-cap4}

\subsection{Equipo de cómputo}
\label{sec:hardware-equipo}

Todos los experimentos del trabajo se ejecutan sobre una única
unidad NVIDIA DGX Spark equipada con el chip Grace Hopper GB10 y
ciento veintiocho gigabytes de memoria unificada compartida entre
CPU y GPU. El sistema operativo es Linux sobre arquitectura ARM de
sesenta y cuatro bits (aarch64), no x86\_64, y la versión de CUDA
instalada es la 13.0. La precisión nativa empleada es BF16
(\emph{brain floating-point} de 16 bits) en la mayor parte de las
operaciones de inferencia y entrenamiento.

Frente al hardware típico de los trabajos publicados (clústeres de
varias GPU A100 o H100 en plataformas x86\_64 con CPU EPYC o Xeon),
esta configuración presenta dos diferencias prácticas relevantes
para el diseño experimental.

\paragraph{Memoria unificada.} El bus PCIe deja de constituir el
cuello de botella habitual entre el host y la GPU, dado que ambos
acceden al mismo banco físico de memoria. La consecuencia operativa
es que pueden mantenerse cargados en GPU modelos del orden de
cincuenta y cinco gigabytes en BF16 sin necesidad de \emph{offloading}
ni de cuantización agresiva. Esta capacidad resulta determinante
para la familia MedGemma evaluada en la
sección~\ref{sec:llm-protocol}, donde la variante de veintisiete mil
millones de parámetros opera en local sin técnicas de compresión
adicionales.

\paragraph{Arquitectura.} La inmensa mayoría de los pesos públicos
empleados en este trabajo se entrenaron y validaron sobre
x86\_64. Diferencias sutiles en la aritmética flotante BF16, en las
bibliotecas BLAS subyacentes y en los \emph{kernels} CUDA
compilados para una u otra arquitectura podrían introducir un sesgo
medible en los resultados. La verificación implícita de que ese
sesgo es despreciable se documenta cuantitativamente en el
capítulo~\ref{cap:resultados}.

\subsection{Entorno de software}
\label{sec:hardware-software}

El entorno de ejecución se compone de Python 3.12.3, PyTorch
2.11.0 y la biblioteca \texttt{timm} 0.9.16 como base. Para las
métricas se emplea \texttt{scikit-learn}; para los modelos
fundacionales publicados, las bibliotecas oficiales de sus autores
(\texttt{open\_clip}, \texttt{transformers}, \texttt{sam2}); para la
adaptación de SAM2.1 se utiliza la biblioteca \texttt{peft} con
LoRA. Para la búsqueda vectorial densa empleada en el módulo de
recuperación visual se utiliza FAISS~\cite{faiss2019}. La variable
de entorno \texttt{WANDB\_MODE} se fija a \texttt{disabled} para
impedir el envío de telemetría a servicios externos, y
\texttt{CUDA\_VISIBLE\_DEVICES} a \texttt{0} para forzar la
ejecución sobre la única GPU presente.

\subsection{Tiempos de ejecución típicos}
\label{sec:hardware-tiempos}

Los tiempos de cálculo característicos sobre este equipo
proporcionan la medida del coste experimental del trabajo:

\begin{itemize}
\item Linear probing completo de un encoder sobre un dataset
estándar: del orden de uno a tres minutos por dataset.
\item Fine-tuning supervisado de PanDerm Large sobre HAM10000 a
cincuenta épocas con \emph{test-time augmentation}: aproximadamente
quince horas.
\item Fine-tuning supervisado de SAM2.1-Large sobre ISIC2018 a
cincuenta épocas con LoRA: aproximadamente veinticuatro horas.
\item Entrenamiento de un Sparse Autoencoder Large (16\,384
características) sobre los embeddings de PanDerm: aproximadamente
cuatro horas.
\item Inferencia sobre el conjunto de evaluación de HAM10000 (1\,232
imágenes) con un encoder fundacional cargado: del orden de
segundos.
\end{itemize}

Estos órdenes de magnitud habilitan la reproducibilidad del trabajo
sobre un único equipo accesible económicamente, sin recurso a
clústeres distribuidos.

\section{Modelos evaluados}
\label{sec:modelos-evaluados}

El trabajo evalúa modelos pertenecientes a cuatro familias: encoders
visuales puros, modelos vision-lenguaje contrastivos, modelo de
segmentación y modelos de lenguaje multimodales generalistas. La
caracterización detallada de cada familia y de sus referentes se
desarrolla en el capítulo~\ref{cap:sota}; esta sección los enumera
con los atributos relevantes para el diseño experimental.

\subsection{Encoders visuales}
\label{sec:modelos-encoders}

\begin{itemize}
\item \textbf{PanDerm Base}~\cite{panderm2025}: ViT-Base/16,
$\sim$~86 millones de parámetros, embedding de salida de 768
dimensiones. Pesos públicos del paper de Yan et al.~(2025).

\item \textbf{PanDerm Large}~\cite{panderm2025}: ViT-Large/16,
$\sim$~307 millones de parámetros, embedding de salida de 1\,024
dimensiones. Referente principal del trabajo.

\item \textbf{DINOv2 ViT-L/14}~\cite{dinov2}: encoder de visión
auto-supervisado de propósito general (Meta AI), $\sim$~300
millones de parámetros. Se incluye como referencia no dermatológica
para contrastar la ventaja de un encoder específico de dominio.

\item \textbf{ConvNeXt-Large}: red convolucional moderna entrenada
de forma supervisada sobre ImageNet-21k, $\sim$~198 millones de
parámetros. Referencia clásica no transformer.

\item \textbf{EfficientNetV2-Large}: red convolucional entrenada de
forma supervisada sobre ImageNet, $\sim$~118 millones de parámetros.
Referencia clásica adicional.
\end{itemize}

\subsection{Modelos vision-lenguaje contrastivos}
\label{sec:modelos-vl}

\begin{itemize}
\item \textbf{DermLIP v1 y v2 (encoder textual PubMedBERT)}~\cite{derm1m}: modelos contrastivos al estilo CLIP con
encoder visual basado en PanDerm Base y encoder textual PubMedBERT
extendido a 256 tokens. La variante DermLIP v2 es la configuración
de mejor rendimiento reportada por los autores en los benchmarks
del propio paper de Derm1M.

\item \textbf{DermLIP original (encoder textual GPT-2/CLIP)}~\cite{derm1m}: variante de la misma familia con encoder visual
ViT-Base y encoder textual GPT-2 con el tokenizador de CLIP (77
tokens). Se incluye como variante de tokenizador para el análisis
de sensibilidad documentado en el capítulo~\ref{cap:resultados}.

\item \textbf{BiomedCLIP}~\cite{biomedclip}: modelo contrastivo
biomédico generalista entrenado sobre $\sim$~15 millones de pares
imagen-texto procedentes de literatura científica.

\item \textbf{SigLIP-Large SO400M}~\cite{siglip2023}: modelo
contrastivo con pérdida sigmoide, embedding visual de 1\,152
dimensiones, $\sim$~878 millones de parámetros. Se incluye en el
trabajo como modelo no específico de dominio para la rama del
ensemble de seguridad descrito en la
sección~\ref{sec:ensemble-protocol}.
\end{itemize}

\subsection{Modelo de segmentación}
\label{sec:modelos-segmentacion}

\textbf{SAM2.1-Large}~\cite{sam2} (\emph{Segment Anything Model
versión 2.1}): modelo de segmentación promptable, $\sim$~227
millones de parámetros, ampliado en este trabajo con adaptadores
LoRA sobre las componentes \texttt{mask\_decoder} y
\texttt{promptable\_encoder}, manteniendo congelado el
\texttt{image\_encoder}. Como referencias se contrastan la versión
\emph{zero-shot} sin ajuste (SAM2 original) y un \emph{Convolutional
AutoEncoder} de segmentación (CAE-seg) entrenado desde cero como
\emph{baseline} clásico.

\subsection{Modelos de lenguaje multimodales generalistas}
\label{sec:modelos-llm}

\begin{itemize}
\item \textbf{GPT-4o} (OpenAI): modelo de lenguaje multimodal de
propósito general, accedido vía API; tamaño no público.

\item \textbf{Gemini 2.5 Pro} (Google DeepMind): modelo multimodal
de propósito general, accedido vía API; tamaño no público.

\item \textbf{MedGemma 4B Vision} (Google): modelo multimodal
abierto, $\sim$~4 mil millones de parámetros, instalado y ejecutado
en local sobre el equipo descrito en la
sección~\ref{sec:hardware-equipo}.

\item \textbf{MedGemma 27B Text} (Google): variante de mayor escala
de la familia MedGemma, $\sim$~27 mil millones de parámetros,
ejecutada en local en BF16 sobre los 128~GB de memoria unificada
disponibles. Se evalúa además una variante ajustada con LoRA en el
mismo conjunto de evaluación.

\item \textbf{BLIP-2}~\cite{blip2_2023}: modelo de visión-lenguaje
general con arquitectura Q-Former, referenciado como contraste
metodológico aunque sin adaptación específica a dermatología en el
presente trabajo.
\end{itemize}

\section{Protocolos experimentales principales}
\label{sec:protocolos-base}

El trabajo aplica cinco protocolos experimentales principales sobre
los modelos descritos. Esta sección formaliza sus hiperparámetros y
decisiones metodológicas; los resultados que cada protocolo genera
se reportan en el capítulo~\ref{cap:resultados}.

\subsection{Linear probing}
\label{sec:lp-protocol}

El protocolo de linear probing mantiene el encoder visual con sus
pesos congelados y entrena un clasificador lineal sobre los
embeddings extraídos. La configuración exacta empleada es la
siguiente:

\begin{itemize}
\item \textbf{Extracción de embeddings.} Cada imagen se preprocesa
con la normalización canónica del modelo (\texttt{ImageNet}
$\mu=[0{,}485, 0{,}456, 0{,}406]$,
$\sigma=[0{,}229, 0{,}224, 0{,}225]$ para PanDerm y DINOv2;
parámetros OpenAI CLIP para los modelos de la familia DermLIP).
El embedding resultante (768 dimensiones en PanDerm Base, 1\,024 en
PanDerm Large, 1\,152 en SigLIP-Large, 1\,024 en DINOv2 ViT-L/14)
se almacena en disco para evitar recomputaciones entre experimentos.

\item \textbf{Clasificador.} Regresión logística multinomial con
algoritmo de optimización L-BFGS (\emph{Limited-memory
Broyden--Fletcher--Goldfarb--Shanno}), de tipo quasi-Newton,
implementado en \texttt{scikit-learn}. Penalización L2 con coste
$C=1{,}0$ por defecto; en algunos experimentos se emplea búsqueda
automática del coste mediante validación cruzada interna.
\texttt{max\_iter}~$=5\,000$ iteraciones.
\texttt{random\_state}~$=42$.

\item \textbf{Splits.} Particiones canónicas de cada dataset
descritas en la tabla~\ref{tab:datasets-resumen} del
capítulo~\ref{cap:datos}.

\item \textbf{Modelos evaluados con LP.} PanDerm Base y Large,
DINOv2, ConvNeXt-Large, EfficientNetV2-Large, los modelos
vision-lenguaje en su rama visual (DermLIP v2, BiomedCLIP,
SigLIP-Large).
\end{itemize}

El protocolo proporciona una medida directa de la calidad de las
representaciones aprendidas durante el preentrenamiento: si las
\emph{features} son adecuadas para la tarea, basta una capa lineal
para alcanzar un buen rendimiento. La elección de regresión
logística sobre alternativas no lineales responde a esa finalidad
diagnóstica.

\subsection{Fine-tuning supervisado}
\label{sec:ft-protocol}

A diferencia del linear probing, el protocolo de fine-tuning
supervisado ajusta los pesos del encoder visual conjuntamente con
los de una cabeza de clasificación específica de la tarea. La
configuración empleada reproduce la descrita en el paper original
de PanDerm~\cite{panderm2025}.

\begin{itemize}
\item \textbf{Optimizador.} AdamW con \emph{weight decay} 0{,}05.
\item \textbf{Tasa de aprendizaje.} $\eta=5\times10^{-4}$ inicial.
\item \textbf{Scheduler.} Cosine annealing con fase inicial de
\emph{warmup} lineal de 10 épocas.
\item \textbf{Layer decay.} 0{,}65 (las capas profundas adaptan a
mayor velocidad que las superficiales).
\item \textbf{Regularización por trayectoria.} \emph{Drop-path}
con probabilidad 0{,}2.
\item \textbf{Suavizado de etiquetas.} \emph{Label smoothing} con
$\epsilon=0{,}1$.
\item \textbf{Augmentaciones de mezcla.} Mixup con $\alpha=0{,}8$
y CutMix con $\alpha=1{,}0$.
\item \textbf{Augmentaciones de imagen.} Recortes aleatorios,
volteo horizontal, \emph{jitter} de color moderado.
\item \textbf{Muestreo.} \emph{Weighted random sampler} con pesos
inversamente proporcionales a la frecuencia de cada clase, para
compensar el desbalance de los datasets con clase dominante muy
representada (notablemente HAM10000).
\item \textbf{Épocas.} 50.
\item \textbf{\emph{Test-time augmentation} (TTA).} En la
evaluación final sobre HAM10000 se promedian los logits sobre 5
augmentaciones deterministas (volteos horizontal y vertical,
rotación, jitter de color) aplicadas a cada imagen del split de
test.
\item \textbf{Precisión numérica.} BF16 en el cuerpo del modelo,
FP32 en el optimizador y en la pérdida (precisión mixta canónica).
\item \textbf{Seed.} \texttt{torch}, \texttt{numpy} y
\texttt{random} fijados a 0.
\end{itemize}

El protocolo se aplica sobre PanDerm Base y Large en los datasets
HAM10000, PAD-UFES-20, DDI y MSKCC, conforme a las cifras
$N_{\mathrm{test}}$ recogidas en la
tabla~\ref{tab:datasets-resumen}.

\subsection{Eficiencia de etiquetas}
\label{sec:label-eff-protocol}

El protocolo de eficiencia de etiquetas mide la dependencia del
rendimiento respecto al volumen de datos de entrenamiento
disponibles, manteniendo congelado el encoder. Sobre HAM10000 y
PAD-UFES-20 se entrenan clasificadores lineales (con los mismos
hiperparámetros que el protocolo de linear probing,
sección~\ref{sec:lp-protocol}) con submuestreos estratificados del
conjunto de entrenamiento al 1\,\%, 5\,\%, 10\,\%, 20\,\%, 50\,\%
y 100\,\%. La estratificación preserva la proporción original de
clases en cada subconjunto. Los submuestreos se generan con
\texttt{random\_state}~$=42$ para garantizar reproducibilidad.

El protocolo se aplica con PanDerm Large como encoder; PanDerm
Base se utiliza como referencia adicional sobre HAM10000.

\subsection{Segmentación binaria de lesión}
\label{sec:seg-protocol}

El protocolo de segmentación binaria de lesión cutánea adapta
SAM2.1-Large al dominio dermatológico mediante fine-tuning con
LoRA. La configuración empleada es la siguiente:

\begin{itemize}
\item \textbf{Componentes ajustadas.} Decodificador de máscaras
(\texttt{mask\_decoder}) y \emph{promptable encoder}.
\item \textbf{Componentes congeladas.} \texttt{image\_encoder}
(Hiera-Large preentrenado), parte mayoritaria de los parámetros del
modelo.
\item \textbf{Rango LoRA.} Configuración estándar de la biblioteca
\texttt{peft} con \emph{rank} $r=8$.
\item \textbf{Optimizador.} AdamW con $\eta=1\times10^{-4}$ y
\emph{weight decay} 0{,}05.
\item \textbf{Pérdida.} Suma ponderada de \emph{cross-entropy} y
\emph{Dice loss} sobre la máscara binaria predicha.
\item \textbf{Tamaño de batch.} 4 imágenes (limitado por el coste
de SAM2.1-Large en BF16).
\item \textbf{Épocas.} Se evalúan dos configuraciones: 50 épocas
(resultado de referencia) y 100 épocas con \emph{early stopping} en
el \emph{best val} (Dice de validación).
\item \textbf{Normalización.} Uniforme con
$\mu=\sigma=0{,}5$ en los tres canales, según la convención de
SAM2.
\item \textbf{Dataset de entrenamiento.}
ISIC2018~\cite{isic2018} (2\,594 imágenes en el split de
evaluación, máscara binaria por imagen).
\item \textbf{Generalización fuera de dominio.} El modelo entrenado
en ISIC2018 se evalúa adicionalmente sobre ISIC2017 y PH2 sin
reentrenamiento, para estimar la capacidad de transferencia a
fuentes con condiciones de captura distintas.
\end{itemize}

Como referencias se contrastan SAM2 \emph{zero-shot} sin
adaptación y un \emph{Convolutional AutoEncoder} de segmentación
(CAE-seg) entrenado desde cero con la misma pérdida y
configuración de optimización.

\subsection{Zero-shot multimodal}
\label{sec:zs-protocol}

El protocolo de clasificación zero-shot multimodal aprovecha el
espacio de representación compartido imagen-texto de los modelos
vision-lenguaje contrastivos. La configuración es la siguiente:

\begin{itemize}
\item \textbf{Codificación de clases.} Cada clase candidata del
dataset se formula como una o varias frases de plantilla. Las
frases se codifican mediante el encoder textual del modelo
correspondiente y, cuando existen varias plantillas para una
misma clase, los vectores resultantes se promedian para obtener un
único vector representativo de la clase.
\item \textbf{Codificación de imágenes.} Cada imagen de
evaluación se codifica mediante el encoder visual con la
normalización canónica del modelo.
\item \textbf{Predicción.} La clase ganadora es la de mayor
similitud coseno con el vector de la imagen.
\item \textbf{Métrica de comparación.} Similitud coseno.
\end{itemize}

Sobre HAM10000 se evalúan tres configuraciones de prompts:

\begin{enumerate}
\item \textbf{Prompts genéricos en inglés} del estilo \emph{``a
dermoscopic image of a melanoma''}, una plantilla por clase.
\item \textbf{Prompts del paper de Derm1M}~\cite{derm1m}: conjunto
A3 de plantillas más extensas, con varios prompts por clase y
promediado posterior.
\item \textbf{Prompts enriquecidos}: descriptores clínicos
adicionales (modalidad, localización anatómica, criterios
diagnósticos visuales) sobre la base del conjunto A3.
\end{enumerate}

El protocolo se aplica con dos tokenizadores distintos:
PubMedBERT extendido a 256 tokens (DermLIP v1/v2) y GPT-2 con el
tokenizador de CLIP de 77 tokens (DermLIP original). Esta dualidad
permite caracterizar la sensibilidad del rendimiento al
tokenizador, cuyo análisis se recoge en el
capítulo~\ref{cap:resultados}.

Sobre HIBA, MSKCC, DDI y Derm7pt clínico se ejecuta una versión
reducida del protocolo, con los prompts genéricos y los Derm1M,
para la formulación binaria melanoma/no-melanoma.

\subsection{Evaluación de modelos de lenguaje multimodales generalistas}
\label{sec:llm-protocol}

El protocolo evalúa modelos de lenguaje multimodales generalistas
sobre HAM10000 como referencia común con el resto de protocolos.
La configuración es la siguiente:

\begin{itemize}
\item \textbf{Prompt clínico estandarizado.} Texto de
aproximadamente trescientas palabras que (i)~describe el contexto
de la tarea (apoyo al diagnóstico dermatológico, no sustitución
del juicio del especialista), (ii)~presenta la imagen como
entrada, (iii)~enumera las siete categorías diagnósticas de
HAM10000 con su descripción breve y (iv)~solicita el diagnóstico
más probable junto con un breve razonamiento estructurado.
\item \textbf{Conjunto de evaluación.} Mismo split de test
oficial de HAM10000 (1\,232 imágenes) que el utilizado en el
linear probing y el fine-tuning, lo que garantiza comparabilidad
numérica entre los tres protocolos.
\item \textbf{Modo.} \emph{Zero-shot} (sin fine-tuning previo).
Para MedGemma 27B se evalúa adicionalmente una variante ajustada
con LoRA (\emph{rank} $r=8$) sobre el split de entrenamiento de
HAM10000.
\item \textbf{Modelos.} GPT-4o (API OpenAI), Gemini 2.5 Pro (API
Google), MedGemma 4B Vision (local), MedGemma 27B Text (local en
BF16 sobre los 128~GB de memoria unificada).
\item \textbf{Métrica.} La etiqueta predicha por el modelo se
extrae del texto generado mediante una rutina de \emph{parsing}
determinista que mapea el nombre del diagnóstico al código
correspondiente de HAM10000; las respuestas no mapeables (menos
del 1\,\% en todos los modelos) se contabilizan como predicción
incorrecta.
\end{itemize}

\section{Protocolos auxiliares}
\label{sec:protocolos-auxiliares}

Junto a los cinco protocolos principales, el trabajo emplea tres
protocolos auxiliares para abordar aspectos específicos del
comportamiento de los modelos (equidad por fototipo, seguridad de
detección de melanoma y consistencia entre datasets ante
solapamiento).

\subsection{Evaluación de equidad por fototipo}
\label{sec:equidad-protocol}

El protocolo de equidad evalúa cinco encoders fundacionales sobre
el dataset Fitzpatrick17k completo (16\,577 imágenes, 114
patologías, 6 fototipos cutáneos I-VI).

\begin{itemize}
\item \textbf{Modelos comparados.} PanDerm Large, PanDerm Base,
DermLIP v2 (rama visual), DINOv2 ViT-L/14, BiomedCLIP.
\item \textbf{Protocolo común.} Linear probing con regresión
logística L-BFGS, penalización L2 con $C=1{,}0$,
\texttt{max\_iter}~$=5\,000$, \texttt{random\_state}~$=42$. Cada
modelo emplea su normalización canónica (ImageNet o CLIP) en la
extracción de embeddings.
\item \textbf{Niveles de evaluación.} Se reportan dos
formulaciones: (i)~clasificación a 114 patologías (grano fino del
dataset), y (ii)~clasificación a 3 clases de malignidad
(benigno / no neoplásico / maligno), que es la formulación
habitualmente empleada en estudios de equidad sobre
Fitzpatrick17k.
\item \textbf{Subgrupos.} Las métricas por fototipo se calculan
sobre los seis subgrupos Fitzpatrick (I, II, III, IV, V, VI). La
métrica de equidad declarada es el \emph{gap} entre los fototipos
extremos:
$\mathrm{Gap}_{\mathrm{I\to VI}} = \mathrm{BAcc}_\mathrm{I} -
\mathrm{BAcc}_\mathrm{VI}$.
\item \textbf{Validación cruzada del script.} El experimento se
re-ejecuta sobre PanDerm Large como verificación, comprobando
coincidencia hasta cuatro decimales con el resultado original.
\end{itemize}

\subsection{Ensemble de seguridad para detección de melanoma}
\label{sec:ensemble-protocol}

El protocolo combina tres clasificadores independientes para
maximizar la sensibilidad en la detección de melanoma sobre el
split de test de HAM10000.

\begin{itemize}
\item \textbf{Clasificadores integrados.}
(i)~PanDerm Large fine-tuned sobre HAM10000 (M1, módulo del
prototipo descrito en el capítulo~\ref{cap:dermapixel}),
(ii)~el clasificador unificado L1/L2/L3 entrenado sobre el corpus
armonizado descrito en el capítulo~\ref{cap:caminos} (M7) mapeado
a las 7 clases de HAM10000, (iii)~SigLIP-Large SO400M con linear
probing sobre el conjunto de entrenamiento de HAM10000.
\item \textbf{Conjunto de evaluación.} Split de test oficial de
HAM10000 (1\,232 imágenes, de las cuales 70 son melanoma).
\item \textbf{Reglas de combinación evaluadas.}
\begin{itemize}
\item \emph{OR top-3 M1+SigLIP}: la imagen se marca como
sospechosa de melanoma si la clase melanoma aparece entre las
tres primeras posiciones del ranking de cualquiera de los dos
clasificadores.
\item \emph{MaxMel M1+SigLIP}: la probabilidad final de melanoma
es el máximo de las probabilidades reportadas por cada
clasificador.
\item \emph{Avg 3 modelos} (sobre la intersección de los 198
ejemplos con predicciones disponibles para los tres modelos):
promedio aritmético de las tres probabilidades.
\end{itemize}
\item \textbf{Métricas reportadas.} Recall de melanoma en top-1 y
top-3, AUROC binarizada melanoma/no-melanoma, accuracy global, y
número de falsos positivos en el escenario \emph{safety screen}.
\end{itemize}

\subsection{Auditoría de leakage frente a Derm1M}
\label{sec:leakage-protocol}

La auditoría de leakage caracteriza el solapamiento entre el split
de test de cada uno de los datasets contrastados y el corpus de
preentrenamiento Derm1M empleado por DermLIP, dado que un
solapamiento no detectado invalidaría la interpretación de los
resultados zero-shot. El procedimiento es el siguiente:

\begin{itemize}
\item \textbf{Hash MD5 por imagen.} Cada imagen del split de test
de los datasets contrastados se reduce a su hash MD5 con la
biblioteca estándar de Python. El mismo procedimiento se aplica a
las imágenes del corpus Derm1M en la fracción a la que el trabajo
tiene acceso.
\item \textbf{Comparación.} La intersección entre los conjuntos
de hashes determina las imágenes presentes simultáneamente en
ambos corpus. La cardinalidad de esa intersección se reporta por
dataset.
\item \textbf{Datasets auditados.} Los doce datasets externos
descritos en el capítulo~\ref{cap:datos}, con la excepción de
SkinCon que no se evalúa sobre clases diagnósticas.
\item \textbf{Resultado anticipado.} Tres datasets presentan
solapamiento documentado en la literatura (Dermnet, HIBA, MSKCC) y
así se reflejan con \emph{caveat} en la
tabla~\ref{tab:datasets-resumen}. La cuantificación exacta del
solapamiento se reporta en el capítulo~\ref{cap:resultados}.
\end{itemize}

\section{Métricas}
\label{sec:metricas-cap4}

Las métricas empleadas a lo largo del trabajo se eligen según el
tipo de tarea: clasificación, segmentación o equidad por
subgrupos. Todas se calculan con la biblioteca \texttt{scikit-learn}
(módulos \texttt{sklearn.metrics}, \texttt{sklearn.model\_selection})
salvo las específicas de segmentación que se implementan en código
propio sobre PyTorch.

\subsection{Métricas de clasificación}
\label{sec:metricas-clasif}

\paragraph{Accuracy.} Proporción de predicciones correctas sobre
el total. Sensible al desbalance: una clase mayoritaria
artificialmente eleva el valor. Útil únicamente como referencia
comparativa entre modelos sobre el mismo dataset y la misma
distribución de clases.

\paragraph{Balanced accuracy (BAcc).} Promedio de la sensibilidad
(\emph{recall}) sobre todas las clases:
$\mathrm{BAcc}=\frac{1}{C}\sum_{c=1}^{C}\mathrm{recall}_c$,
donde $C$ es el número de clases. Penaliza a los modelos que
concentran su rendimiento en las clases mayoritarias y constituye
la métrica principal del trabajo en escenarios con desbalance.

\paragraph{F1 ponderada (W-F1).} Media armónica de precisión y
\emph{recall} ponderada por la frecuencia de cada clase. Refleja
el rendimiento agregado conservando sensibilidad al desbalance.

\paragraph{Área bajo la curva ROC (AUROC).} Para problemas
binarios, el área bajo la curva que relaciona la tasa de
verdaderos positivos con la de falsos positivos al barrer todos
los umbrales. Para problemas multiclase, se calcula con la
estrategia \emph{one-vs-rest} de \texttt{scikit-learn}
(\texttt{multi\_class='ovr'}) y se reporta el valor medio. Es la
métrica más robusta al desbalance y la más comparable entre
trabajos publicados.

\paragraph{Coeficiente kappa de Cohen.} Mide la concordancia entre
predicción y etiqueta corrigiendo por la concordancia esperada por
azar. Útil para datasets multiclase muy desbalanceados.

\paragraph{Recall@$k$.} Para clasificación con $C$ clases,
\emph{recall@$k$} mide si la etiqueta verdadera aparece entre las
$k$ primeras posiciones del ranking de probabilidades del modelo.
El trabajo lo reporta específicamente para la clase melanoma
sobre HAM10000, dado que en un escenario de tamizaje clínico es
preferible que una lesión maligna aparezca entre las primeras
opciones, aunque no sea la predicción principal, a que pase
desapercibida.

\subsection{Métricas de segmentación}
\label{sec:metricas-segm}

\paragraph{Coeficiente Dice (DSC).}
$\mathrm{DSC} = \frac{2\,|P \cap G|}{|P| + |G|}$,
donde $P$ es la máscara predicha y $G$ la máscara de referencia.
Métrica canónica para segmentación binaria de lesión.

\paragraph{Intersección sobre unión (IoU).}
$\mathrm{IoU} = \frac{|P \cap G|}{|P \cup G|}$,
también conocida como índice de Jaccard. Complementaria al Dice.

\subsection{Métricas de equidad}
\label{sec:metricas-equidad}

\paragraph{BAcc por subgrupo Fitzpatrick.} El protocolo de equidad
(sección~\ref{sec:equidad-protocol}) reporta la BAcc separadamente
para cada uno de los seis fototipos cutáneos I-VI, sobre la
formulación de 3 clases de malignidad.

\paragraph{Gap I-VI.} Diferencia entre la BAcc del fototipo I y la
del fototipo VI. Valores próximos a cero indican que el rendimiento
es comparable entre fototipos extremos. Esta cifra se interpreta
junto con la BAcc global: un \emph{gap} pequeño asociado a una
BAcc global también pequeña reflejaría un modelo uniformemente
inferior, no un modelo equitativo en sentido sustantivo.

\subsection{Métricas operativas}
\label{sec:metricas-operativas}

Junto a las métricas de rendimiento se reportan dos métricas
operativas relevantes para el despliegue:

\paragraph{Latencia por consulta.} Tiempo desde la entrada de la
imagen hasta la salida del modelo (inferencia más posproceso),
medido en milisegundos sobre el equipo descrito en
la~\ref{sec:hardware-equipo}.

\paragraph{Huella de memoria.} VRAM ocupada por el modelo cargado
en GPU durante la inferencia, expresada en gigabytes en BF16.

\section{Trazabilidad y reproducibilidad}
\label{sec:reproducibilidad}

El diseño experimental incorpora las prácticas habituales para
garantizar la reproducibilidad de los resultados reportados en el
capítulo~\ref{cap:resultados}.

\paragraph{Semillas aleatorias.} Las semillas de
\texttt{torch}, \texttt{numpy} y \texttt{random} se fijan a
$0$ al inicio de cada experimento de fine-tuning supervisado y de
segmentación. Para los experimentos basados en \texttt{scikit-learn}
(linear probing, eficiencia de etiquetas, equidad por fototipo) la
semilla declarada en el parámetro \texttt{random\_state} es
$42$.

\paragraph{Versionado del software.} El entorno de ejecución
queda fijado en Python 3.12.3, PyTorch 2.11.0, \texttt{timm}
0.9.16 y CUDA 13.0. Cualquier desviación respecto a estas
versiones se documenta explícitamente en el experimento
correspondiente.

\paragraph{Splits y vocabularios.} Todos los protocolos emplean
los splits canónicos descritos en la
tabla~\ref{tab:datasets-resumen} del capítulo~\ref{cap:datos}. Los
mapeos de etiquetas nativas a la ontología jerárquica L1/L2/L3 se
mantienen fijos a través de los experimentos y se documentan en el
anexo~\ref{anexo:tareas}.

\paragraph{Aislamiento del entorno de ejecución.} La variable
\texttt{WANDB\_MODE} se fija a \texttt{disabled} para impedir el
envío de telemetría a servicios externos durante el entrenamiento.
\texttt{CUDA\_VISIBLE\_DEVICES} se fija a \texttt{0} para forzar la
ejecución sobre la única GPU presente en el equipo.

\paragraph{Almacenamiento de resultados.} Los resultados numéricos
de cada experimento se vuelcan a ficheros CSV o JSON con esquema
estable, lo que permite la reconstrucción a posteriori de cualquier
tabla o figura del capítulo~\ref{cap:resultados} sin necesidad de
reejecución. El listado de tareas reproducibles, sus comandos de
ejecución, ficheros de salida y \emph{checkpoints} de modelo se
documenta en el anexo~\ref{anexo:tareas}.

\paragraph{Trazabilidad inter-experimentos.} Cuando un experimento
emplea como entrada el resultado de otro (por ejemplo, los
embeddings extraídos en la fase de linear probing sirven de base a
los modelos vision-lenguaje contrastivos en zero-shot), la
identidad de la entrada se preserva mediante hash MD5 sobre el
fichero correspondiente. Esto permite verificar a posteriori que
los experimentos encadenados parten del mismo material y son
mutuamente consistentes.

Con el hardware, el conjunto de modelos, los siete protocolos
formalizados y las métricas descritas, queda definido el marco
experimental sobre el que se reportan los resultados del
capítulo~\ref{cap:resultados} y se discuten en el mismo. El
capítulo~\ref{cap:caminos} retoma esos elementos para enmarcar las
líneas de continuación no abordadas en el periodo cubierto por
este documento.
